
\documentclass[12pt]{article}
\usepackage[margin=0.75in]{geometry}
\usepackage{fontspec}
\setmainfont{Latin Modern Roman}
\usepackage{microtype}
\usepackage{multicol}
\usepackage{fancyhdr}
\usepackage{titlesec}
\usepackage{enumitem}
\usepackage{xcolor}
\usepackage[hidelinks]{hyperref}
\setlength{\columnsep}{0.28in}
\setlength{\parindent}{1.1em}
\setlength{\parskip}{0.35em}
\setlength{\headheight}{15pt}
\titleformat{\section}{\large\bfseries\scshape}{}{0pt}{}
\titleformat{\subsection}{\normalsize\bfseries}{}{0pt}{}
\pagestyle{fancy}
\fancyhf{}
\fancyfoot[C]{\thepage}
\renewcommand{\headrulewidth}{0.4pt}
\newcommand{\focusbox}[1]{\par\vspace{0.4em}\noindent\begingroup\setlength{\fboxsep}{6pt}\colorbox{gray!12}{\begin{minipage}{\dimexpr\linewidth-2\fboxsep\relax}#1\end{minipage}}\endgroup\par\vspace{0.5em}}
\newcommand{\studentline}{\vspace{0.35em}\hrule\vspace{0.65em}}

\fancyhead[L]{Whately: Logic Is Not Magic}
\fancyhead[R]{Student Reading}
\begin{document}
\begin{center}
{\LARGE\bfseries Logic Is Not Magic}\par\vspace{0.25em}
{\large\scshape What Logic Can and Cannot Do}\par\vspace{0.5em}
{Adapted and modernized from Richard Whately, \textit{Elements of Logic}}\par\vspace{0.6em}
\rule{0.82\textwidth}{0.4pt}
\end{center}

\section*{Central Question}
What can logic actually do for a thinker, and what should we not expect it to do?

\section*{Vocabulary Preview}
\begin{multicols}{2}
\begin{itemize}[leftmargin=*, itemsep=0.18em]
\item \textbf{Logic}: the study of the structure of reasoning.
\item \textbf{Reasoning}: moving from certain claims to another claim that follows from them.
\item \textbf{Argument}: a set of reasons offered to support a conclusion.
\item \textbf{Premise}: a claim used as support in an argument.
\item \textbf{Conclusion}: the claim an argument is trying to prove.
\item \textbf{Deduction}: reasoning in which a conclusion is drawn from premises.
\item \textbf{Fallacy}: reasoning that appears persuasive but does not actually prove the point.
\item \textbf{Science}: organized knowledge of principles.
\item \textbf{Art}: the practical use of knowledge to do something well.
\item \textbf{Panacea}: a supposed cure for every problem.
\end{itemize}
\end{multicols}

\section*{Introduction}
Richard Whately was a nineteenth-century writer, teacher, and Archbishop of Dublin. His book \textit{Elements of Logic} tried to rescue logic from two opposite mistakes. Some people treated logic as useless word games. Others promised too much for it, as if logic could discover every truth and make every person wise. Whately's answer is more careful: logic cannot do everything, but it can train us to see whether an argument really works.

\focusbox{\textbf{Source note:} This reading is adapted and modernized from Whately's introduction to \textit{Elements of Logic}. It preserves his main sequence of ideas while putting the language into modern student-facing prose.}

\begin{multicols}{2}
\section*{Reading}
Logic, in the broadest useful sense of the word, may be considered both a \textit{science} and an \textit{art} of reasoning.

As a science, logic studies the principles by which arguments are made. It asks: What is happening when the mind moves from one claim to another? What makes a conclusion follow from its premises? What kinds of mistakes make an argument only appear to prove its point?

As an art, logic gives practical rules for guarding against false deductions. A person who has studied logic should be better prepared to notice when an argument has skipped a step, changed the meaning of a word, assumed what it needed to prove, or proved something different from the question at hand.

But this does not mean that logic is magic. Logic does not give sight to the blind. It does not supply facts where no facts have been learned. It does not make a foolish person wise merely by giving him technical terms. It does not discover the truth of every subject by itself. If someone expects logic to do all these things, he will be disappointed, and then he may blame logic unfairly.

This is one reason logic has often been despised. Some of its defenders praised it as if it were a cure for every intellectual disease. They spoke as though logic could teach the whole use of reason, settle every question, and discover truth in every field of study. When logic failed to perform these impossible tasks, critics treated the whole subject as empty.

That is unfair. A tool should be judged by its proper use. We do not blame arithmetic because it cannot teach history. We do not blame a map because it cannot walk the road for us. We do not blame the science of optics because it cannot give eyes to a blind man. In the same way, we should not blame logic because it cannot do the work of observation, memory, experience, imagination, or moral judgment.

Logic has a narrower but still important office. It examines reasoning as reasoning. It helps us see whether the conclusion follows from what has been granted. If the premises are false, logic alone may not tell us that; we may need history, science, testimony, experience, or careful observation. But if the reasoning is bad, logic helps expose the fault.

This is why logic matters for students, readers, and citizens. We are constantly surrounded by claims. Some are true, some false, some confused, and some only half-proved. People argue in essays, speeches, advertisements, conversations, courtrooms, sermons, and political debates. Many arguments sound strong because they are spoken confidently or because they appeal to feelings. But confidence is not proof. Emotion is not proof. Repetition is not proof.

Logic asks a quieter question: What follows from what?

If a person says, "This policy is popular, therefore it is just," logic asks whether popularity proves justice. If a person says, "This character failed, therefore he deserved to fail," logic asks what hidden premise connects failure with desert. If a person says, "Everyone knows this is true," logic asks whether common opinion is being used as evidence or merely as pressure.

A student who learns logic is not learning a trick for winning arguments. He is learning a habit of intellectual honesty. He learns to slow down before agreeing. He learns to ask what has actually been proved. He learns to distinguish a real argument from a sentence that only sounds like one.

Logic, then, is modest and powerful at the same time. It is modest because it does not pretend to replace knowledge, wisdom, or virtue. It is powerful because it gives the mind discipline. It trains us to notice the structure beneath words. It helps us see whether a conclusion has been earned.

A person may reason before studying logic, just as a person may speak before studying grammar. But grammar can make speech clearer, and logic can make reasoning more careful. Logic does not create reason; it studies reason so that we may use it better.

\section*{Reading Focus}
\begin{enumerate}[leftmargin=*]
\item In one sentence, state Whately's main claim about what logic is.
\item What are two things logic \textit{can} do?
\item What are two things logic \textit{cannot} do?
\item Why does Whately think people have unfairly blamed logic?
\item Choose one analogy from the reading: map, arithmetic, optics, grammar, or another. Explain how it helps define logic.
\end{enumerate}

\section*{Handout Summary}
Write 3--5 sentences explaining Whately's view of logic. Your summary must include the words \textit{reasoning}, \textit{argument}, and \textit{proof}.
\end{multicols}
\end{document}
